Let $\mathcal{A}$ denote the set of available arms. At time step $k$, the
algorithm selects an arm $A_k\in\mathcal{A}$ and receives a random reward
$R_k$. The true action value of arm $a$ is
\[
    q^\ast(a)
    =
    \mathbb{E}\!\left[R_k \mid A_k=a\right],
\]
which is the expected reward obtained when arm $a$ is selected.

Let
\[
    a^\ast \in \arg\max_{a\in\mathcal{A}} q^\ast(a)
\]
be an optimal arm. The best achievable expected reward per time step is then
\[
    R^\ast
    =
    q^\ast(a^\ast)
    =
    \max_{a\in\mathcal{A}} q^\ast(a).
\]
Thus, $R^\ast$ is a deterministic expected reward, whereas $R_k$ is the
random reward actually observed at time step $k$.

The cumulative expected regret after $t$ time steps is defined as
\begin{equation}
    \operatorname{Regret}_t
    =
    \mathbb{E}\!\left[
        \sum_{k=1}^{t}
        \left(R^\ast-R_k\right)
    \right].
    \label{eq:q1-regret}
\end{equation}
The expectation in \eqref{eq:q1-regret} accounts for the randomness in the
observed rewards and, where applicable, the randomisation of the bandit
policy. In particular, if arm $a$ is selected at time $k$, its expected
one-step regret is
\[
    R^\ast-q^\ast(a)\geq 0.
\]
Hence, regret measures the expected cumulative loss in reward caused by not
knowing the optimal arm in advance.

A bandit algorithm is said to have \emph{sublinear regret} if
\begin{equation}
    \lim_{t\to\infty}
    \frac{\operatorname{Regret}_t}{t}
    =
    0,
    \label{eq:q1-sublinear-regret}
\end{equation}
or equivalently,
\[
    \operatorname{Regret}_t=o(t).
\]

Using \eqref{eq:q1-regret},
\begin{align}
    \frac{\operatorname{Regret}_t}{t}
    &=
    \frac{1}{t}
    \mathbb{E}\!\left[
        \sum_{k=1}^{t}
        \left(R^\ast-R_k\right)
    \right]
    \notag\\
    &=
    R^\ast
    -
    \mathbb{E}\!\left[
        \frac{1}{t}\sum_{k=1}^{t}R_k
    \right].
    \label{eq:q1-average-regret}
\end{align}
Therefore, by \eqref{eq:q1-sublinear-regret} and
\eqref{eq:q1-average-regret}, sublinear regret implies
\[
    \lim_{t\to\infty}
    \mathbb{E}\!\left[
        \frac{1}{t}\sum_{k=1}^{t}R_k
    \right]
    =
    R^\ast.
\]

Thus, the significance of sublinear regret is that the average expected
reward of the learning algorithm converges to the best achievable expected
reward. Although the algorithm incurs regret while exploring and learning
which arm is optimal, the regret per time step vanishes asymptotically.